\section{Offline results: how much verification can the agent's internals replace?}\label{sec:results}

This section answers RQ1--RQ3 on the sealed evaluation sets, with all screens frozen beforehand. Unless stated
otherwise, figures refer to the tolerance $X=10\%$, the path-planning fresh certification set, and the CSTR
certification partition. Tables~\ref{tab:skip-fsm} and~\ref{tab:skip-cstr} list the share of validator calls
skipped by every signal set and model. Figure~\ref{fig:composition} decomposes the skipped calls into unchecked
accepts and unchecked rejects.

\begin{table*}[!tbp]
\centering
\small
\caption{Path planning: share of validator calls skipped (\%) per signal on the fresh certification set (tolerance $X=10\%$). Bold: failure rate among unchecked accepts certified $\le 10\%$ (one-sided Clopper--Pearson, Bonferroni over signals). $^{+}$/$^{-}$: significantly more/fewer calls skipped than observables (paired bootstrap 95\% CI).}
\label{tab:skip-fsm}
\resizebox{\textwidth}{!}{%
\begin{tabular}{lrrrrr}
\toprule
Signal & Qwen-1.5B & Qwen-3B & Qwen-7B & Llama-3B & SmolLM-1.7B \\
\midrule
Observables & \textbf{40} & \textbf{21} & \textbf{16} & 43 & 85 \\
\midrule
Obs. + grounding & \textbf{46}$^{+}$ & \textbf{21} & 22$^{+}$ & \textbf{59}$^{+}$ & 85 \\
Obs. + all internals & \textbf{37}$^{-}$ & \textbf{32}$^{+}$ & 15 & \textbf{47}$^{+}$ & 79$^{-}$ \\
Obs. + internals + grounding & \textbf{40} & \textbf{34}$^{+}$ & \textbf{15} & \textbf{58}$^{+}$ & 91$^{+}$ \\
\midrule
Region grounding & \textbf{31}$^{-}$ & \textbf{14}$^{-}$ & \textbf{3}$^{-}$ & \textbf{31}$^{-}$ & 80$^{-}$ \\
Attention (regions) & \textbf{35}$^{-}$ & \textbf{21} & 13$^{-}$ & 17$^{-}$ & 75$^{-}$ \\
Hidden states & 22$^{-}$ & \textbf{24}$^{+}$ & 10$^{-}$ & 27$^{-}$ & 74$^{-}$ \\
All internals & \textbf{25}$^{-}$ & \textbf{25}$^{+}$ & \textbf{9}$^{-}$ & \textbf{31}$^{-}$ & 76$^{-}$ \\
All internals + grounding & \textbf{38}$^{-}$ & \textbf{23}$^{+}$ & \textbf{11}$^{-}$ & \textbf{36}$^{-}$ & 79$^{-}$ \\
Token confidence & 13$^{-}$ & 18$^{-}$ & 0$^{-}$ & 14$^{-}$ & 56$^{-}$ \\
\bottomrule
\end{tabular}}
\end{table*}

\begin{table*}[!tbp]
\centering
\small
\caption{CSTR: share of validator calls skipped (\%) per signal on the certification partition ($X=10\%$); notation as in Table~\ref{tab:skip-fsm}. $^\dagger$Qwen-3B: exploratory re-analysis under the bound-based rule (its sealed sets had been used once before the rule change).}
\label{tab:skip-cstr}
\resizebox{\textwidth}{!}{%
\begin{tabular}{lrrrrr}
\toprule
Signal & Qwen-1.5B & Qwen-3B$^\dagger$ & Qwen-7B & Llama-3B & SmolLM-1.7B \\
\midrule
Observables & 63 & 41 & 86 & 71 & 61 \\
\midrule
Obs. + grounding & 60$^{-}$ & 39 & 86 & 75$^{+}$ & 65$^{+}$ \\
Obs. + all internals & 61 & 51 & 83 & 72 & 66$^{+}$ \\
Obs. + internals + grounding & 61 & 50 & 80$^{-}$ & 74 & 66$^{+}$ \\
\midrule
Region grounding & 27$^{-}$ & 10 & 47$^{-}$ & 38$^{-}$ & 24$^{-}$ \\
Attention (regions) & 30$^{-}$ & 15 & 49$^{-}$ & 48$^{-}$ & 31$^{-}$ \\
Hidden states & 29$^{-}$ & 13 & 50$^{-}$ & 48$^{-}$ & 45$^{-}$ \\
All internals & 29$^{-}$ & 15 & 47$^{-}$ & 53$^{-}$ & 40$^{-}$ \\
All internals + grounding & 28$^{-}$ & 12 & 48$^{-}$ & 53$^{-}$ & 41$^{-}$ \\
Token confidence & 4$^{-}$ & 0 & 17$^{-}$ & 34$^{-}$ & 22$^{-}$ \\
\bottomrule
\end{tabular}}
\end{table*}


\subsection{RQ1: internals alone, without any model of the plant}\label{sec:res-rq1}

\paragraph{Path planning} The internals-only screen (all internals with grounding) skipped 38.0\%, 22.6\%, 10.8\%,
35.8\% and 78.8\% of validator calls for Qwen2.5-1.5B, Qwen2.5-3B, Qwen2.5-7B, Llama-3.2-3B and SmolLM2-1.7B. For
the first four models the failure rate among its unchecked accepts was \emph{certified} below 10\%: 8 failures
among 240 accepts, 21 among 483, 12 among 303 and 13 among 371, with Bonferroni-corrected upper bounds of 6.8--7.8\%.
For Llama-3.2-3B the internals-only screen was certified even though the observable screen was not (upper bound
6.8\% against 12.4\%). For SmolLM2, whose proposals fail 83\% of the time, the internals-only screen accepted
nothing and skipped its calls entirely by rejection. 95.2\% of the proposals it rejected did fail. Against the
observable screen, internals alone skipped significantly more calls for Qwen2.5-3B (+2.0 points [0.4, 3.7]) and
fewer for the other four models (between $-2.4$ and $-7.2$ points).

\paragraph{CSTR} Without any plant reading or proposed setpoint, the internals-only screen skipped 28.3\%, 12.2\%,
47.8\%, 52.7\% and 40.5\% of validator calls (same model order). Almost all of these were \emph{rejections}
(Fig.~\ref{fig:composition}b). The screen accepted no proposal for four models. For SmolLM2 it accepted 39
proposals, of which 2 failed. Its rejections were precise: 95.1--97.6\% of the proposals it returned to the agent
without simulation would indeed have failed validation. The observable screens skipped more calls (41--86\%),
including unchecked accepts for three models, and had similar rejection precision (92.2--97.1\%). Single internal
families were less reliable than their combination. Hidden states alone, for example, accepted 65 SmolLM2
proposals, of which 9 failed. Token confidence was the weakest signal in both domains: it skipped 0--56\% of calls
in path planning and 0--34\% in the CSTR and was never certified.

\begin{figure}[!htbp]
\centering
% How each screen skips validator calls: unchecked accepts (dark) and unchecked rejects (light), as % of proposals.
% Data: data/composition_{fsm,cstr}.dat (scripts/50).
\begin{tikzpicture}
\begin{groupplot}[group style={group size=2 by 1, horizontal sep=12mm},
  paperaxis, modelticks, width=0.47\textwidth, height=0.42\textwidth, ymin=0, ymax=100,
  ylabel={proposals decided without validation (\%)}, ytick={0,20,40,60,80,100}]
\nextgroupplot[title={\footnotesize (a) Path planning, fresh certification set}]
\foreach \k/\c/\s in {obs/cObs/-6.5pt, int/cInt/0pt, comb/cComb/6.5pt} {
  \edef\tmp{\noexpand\addplot[ybar, bar width=5.5pt, bar shift=\s, draw=\c, fill=\c!30, forget plot]
    table[x=x, y expr=\noexpand\thisrow{\k_acc}+\noexpand\thisrow{\k_rej}] {data/composition_fsm.dat};
  \noexpand\addplot[ybar, bar width=5.5pt, bar shift=\s, draw=\c, fill=\c, forget plot]
    table[x=x, y=\k_acc] {data/composition_fsm.dat};}\tmp
}
\nextgroupplot[title={\footnotesize (b) CSTR, certification partition}, ylabel={},
  legend columns=3, legend cell align=left, legend style={/tikz/every even column/.append style={column sep=8pt}},
  legend to name=complegend]
\addlegendimage{area legend, draw=cObs, fill=cObs}\addlegendentry{observables}
\addlegendimage{area legend, draw=cInt, fill=cInt}\addlegendentry{internals only (no plant model)}
\addlegendimage{area legend, draw=cComb, fill=cComb}\addlegendentry{observables + internals + grounding}
\foreach \k/\c/\s in {obs/cObs/-6.5pt, int/cInt/0pt, comb/cComb/6.5pt} {
  \edef\tmp{\noexpand\addplot[ybar, bar width=5.5pt, bar shift=\s, draw=\c, fill=\c!30, forget plot]
    table[x=x, y expr=\noexpand\thisrow{\k_acc}+\noexpand\thisrow{\k_rej}] {data/composition_cstr.dat};
  \noexpand\addplot[ybar, bar width=5.5pt, bar shift=\s, draw=\c, fill=\c, forget plot]
    table[x=x, y=\k_acc] {data/composition_cstr.dat};}\tmp
}
\end{groupplot}
\node[anchor=north] at ($(group c1r1.south)!0.5!(group c2r1.south)+(0,-1.05)$) {\pgfplotslegendfromname{complegend}};
\end{tikzpicture}

\caption{How each screen avoids the validator ($X=10\%$). Full bar: proposals decided without simulation. Dark
part: unchecked accepts (executed without simulation). Light part: unchecked rejects (returned to the agent
without simulation). In the CSTR the internals-only screen, which uses no plant reading and no proposed
setpoint, skips 12--53\% of calls almost entirely by rejection. Qwen-3B in the CSTR is an exploratory
re-analysis.}
\label{fig:composition}
\end{figure}

\paragraph{Answer to RQ1} The agent's internals are, on their own, a usable replacement for part of the
verification. In a retrieval-heavy task they support certified unchecked execution. In a physics-heavy task
they support high-precision unchecked rejection but not certifiable acceptance. One explanation consistent
with the rest of our results (Section~\ref{sec:res-rq3}) is that internal signals detect a subset of failures:
those accompanied by an atypical internal state, such as low confidence, diffuse or misplaced attention, or an
unusual representation. Detecting a subset of failures suffices to reject with high precision. Accepting safely
requires ruling out \emph{all} failure modes, including confident, well-grounded reasoning errors, and these
leave no internal trace.

\subsection{RQ2: do internals add to observable features?}\label{sec:res-rq2}

\paragraph{Path planning} Adding internal signals to observables significantly increased the share of calls
skipped for every model, with at least one combination (Fig.~\ref{fig:delta}a). The full combination added 14.0
points [12.3, 15.6] for Qwen2.5-3B, 15.0 [13.2, 16.9] for Llama-3.2-3B and 6.2 [4.9, 7.5] for SmolLM2. Grounding
alone added 5.1 points [3.6, 6.7] for Qwen2.5-1.5B, 15.5 [13.8, 17.2] for Llama-3.2-3B and 6.5 [5.2, 7.7] for
Qwen2.5-7B. Most of these screens were certified. The exceptions are informative. For Qwen2.5-7B the extra calls
skipped with grounding came with an accept set that failed certification (41 failures among 513 accepts, upper
bound 11.6\%), and its certified combinations did not beat observables. For SmolLM2 no screen was certifiable,
including the observable one (15 failures among 166 accepts). The gains came from both decisions: the full
combination raised unchecked accepts from 11.2\% to 20.2\% of proposals for Qwen2.5-3B, and unchecked rejects from
31.9\% to 40.3\% for Llama-3.2-3B (Fig.~\ref{fig:composition}a).

\paragraph{CSTR} The picture differs (Fig.~\ref{fig:delta}b). Internal signals added significantly for two
models. For SmolLM2 the gains were +3.8 points [1.4, 6.3] with grounding, +4.8 [1.4, 8.5] with all internals and
+4.6 [1.0, 8.3] with both. For Llama-3.2-3B grounding added +4.2 points [0.8, 7.8], and all internals did not
change the share skipped but halved the good proposals rejected unchecked (from 16 to 8 of 79; difference
$-8$ [$-15$, $-2$]). For two models adding internals \emph{reduced} the calls skipped: Qwen2.5-1.5B with
grounding ($-2.8$ [$-5.5$, $-0.4$]) and Qwen2.5-7B with the full combination ($-5.6$ [$-8.6$, $-2.6$]). For
Qwen2.5-3B, whose analysis is exploratory, the combination raised the share skipped from 41.2\% to 49.8\%. It
also made unchecked accepts possible at all: 43 accepts, of which 5 failed.

No CSTR screen was certified at $X=10\%$. This is the power limit of Eq.~\eqref{eq:mmin}, stated in advance. The
largest accept set of any CSTR screen held 128 proposals, and the observed failure rates among unchecked accepts
were 3.7--11.6\% for the observable screens that accepted anything. For two models (Qwen2.5-1.5B, 11.6\%; SmolLM2,
10.9\%) the observed rate exceeded the tolerance. The certificate correctly refused these screens. This is
the safeguard working as intended: a threshold chosen on development data does not by itself guarantee the
tolerance on new data.

\begin{figure}[!htbp]
\centering
% Change in validator calls skipped when internal signals are added to observables (X = 10%), paired bootstrap 95% CI.
% Data: data/delta_{fsm,cstr}.dat (scripts/50).
\begin{tikzpicture}
\begin{groupplot}[group style={group size=2 by 1, horizontal sep=12mm},
  paperaxis, modelticks, width=0.47\textwidth, height=0.4\textwidth,
  ylabel={$\Delta$ calls skipped vs.\ observables (points)}]
\nextgroupplot[title={\footnotesize (a) Path planning}, ymin=-8, ymax=19]
\addplot[black!50, forget plot] coordinates {(0.4,0) (5.6,0)};
\addplot[only marks, mark=*, mark size=1.8pt, cGrd, error bars/.cd, y dir=both, y explicit]
  table[x expr=\thisrow{x}-0.22, y=grd, y error plus=grd_hi, y error minus=grd_lo] {data/delta_fsm.dat};
\addplot[only marks, mark=square*, mark size=1.7pt, cInt, error bars/.cd, y dir=both, y explicit]
  table[x expr=\thisrow{x}, y=int, y error plus=int_hi, y error minus=int_lo] {data/delta_fsm.dat};
\addplot[only marks, mark=triangle*, mark size=2.2pt, cComb, error bars/.cd, y dir=both, y explicit]
  table[x expr=\thisrow{x}+0.22, y=both, y error plus=both_hi, y error minus=both_lo] {data/delta_fsm.dat};
\nextgroupplot[title={\footnotesize (b) CSTR}, ymin=-10, ymax=10, ylabel={},
  legend columns=3, legend cell align=left, legend style={/tikz/every even column/.append style={column sep=8pt}},
  legend to name=deltalegend]
\addplot[black!50, forget plot] coordinates {(0.4,0) (5.6,0)};
\addplot[only marks, mark=*, mark size=1.8pt, cGrd, error bars/.cd, y dir=both, y explicit]
  table[x expr=\thisrow{x}-0.22, y=grd, y error plus=grd_hi, y error minus=grd_lo] {data/delta_cstr.dat};
\addlegendentry{+ grounding}
\addplot[only marks, mark=square*, mark size=1.7pt, cInt, error bars/.cd, y dir=both, y explicit]
  table[x expr=\thisrow{x}, y=int, y error plus=int_hi, y error minus=int_lo] {data/delta_cstr.dat};
\addlegendentry{+ all internals}
\addplot[only marks, mark=triangle*, mark size=2.2pt, cComb, error bars/.cd, y dir=both, y explicit]
  table[x expr=\thisrow{x}+0.22, y=both, y error plus=both_hi, y error minus=both_lo] {data/delta_cstr.dat};
\addlegendentry{+ internals + grounding}
\node[font=\scriptsize, text=black!60, align=center] at (axis cs:2,6.5) {exploratory\\[-1pt](no CI)};
\end{groupplot}
\node[anchor=north] at ($(group c1r1.south)!0.5!(group c2r1.south)+(0,-1.05)$) {\pgfplotslegendfromname{deltalegend}};
\end{tikzpicture}

\caption{Change in validator calls skipped when grounding, all internal signals, or both are added to the
observable screen ($X=10\%$; paired bootstrap 95\% CI). An interval above zero means the internals let the agent
skip significantly more simulations at the same tolerance. Qwen-3B in the CSTR is an exploratory re-analysis
without intervals.}
\label{fig:delta}
\end{figure}

\paragraph{Answer to RQ2} Internals complement observables strongly in path planning (all five models) and
weakly in the CSTR (two of four confirmatory models, with losses for two). Where internals do not add, the
observable state together with the proposed action already carries the information. In the CSTR on development
data, the observable probe reached an AUROC of 0.93 in the stratum where doing nothing would have passed, so
whether a proposed change breaks a recovering plant is largely readable from the change and the plant state.

\subsection{RQ3: why internals help in one domain and not in the other}\label{sec:res-rq3}

The region-based grounding signal allows a direct test of the mechanism. If failures in a domain are retrieval
failures, low attention on the deciding region should predict failure. If they are reasoning failures over
correctly retrieved information, it should not. The expected contrast between the domains was written into the
CSTR grounding pre-registration before any CSTR grounding value was computed.

\paragraph{Path planning: retrieval failures} Low grounding predicted failure strongly for every model
(Fig.~\ref{fig:mechanism}): the AUROC of $-g_{\min}$ for an invalid answer was 0.72--0.88. The effect is visible at
the level of individual steps. An illegal step received markedly less attention on its source node's adjacency
line than a legal one (step-level AUROC 0.64--0.81 on development data and 0.67--0.81 on test data, four models).
For Llama-3.2-3B the deciding line was the most-attended line for 51\% of legal steps but only 20\% of illegal
ones. As a feature, grounding raised the AUROC of the observable probe by 0.03--0.05 on the test partitions
(four models, all intervals above zero), whereas the coarse internal features did so for only one model.

\paragraph{CSTR: reasoning failures} The same signal was absent for three models (AUROC 0.48--0.52) and weak to
moderate for the two smallest (0.66 for SmolLM2, 0.72 for Qwen2.5-1.5B). A pre-registered test for Qwen2.5-3B
makes the contrast concrete. 339 of 3954 development proposals raised the feed while cooling was saturated and the
reactor was hot. The physics stated in the prompt rules out this move, and only 5.9\% of these proposals passed.
The hypothesis was that such proposals would attend \emph{less} to the physics section. They attended slightly
\emph{more}: the AUROC in the predicted direction was 0.31, that is, 0.69 in the opposite direction. The model reads
the relevant physics and still moves the feed the wrong way.

\begin{figure}[!htbp]
\centering
% Does low attention on the deciding prompt region predict failure? AUROC per model and domain.
% Data: data/mechanism.dat (scripts/50).
\begin{tikzpicture}
\begin{axis}[paperaxis, modelticks, width=0.62\textwidth, height=0.4\textwidth, ymin=0.4, ymax=1.0,
  ylabel={AUROC of $-g_{\min}$ for failure}, ytick={0.4,0.5,0.6,0.7,0.8,0.9},
  legend columns=2, legend style={at={(0.5,0.98)}, anchor=north, /tikz/every even column/.append style={column sep=8pt}},
  legend cell align=left]
\addplot[black!60, densely dashed, forget plot] coordinates {(0.4,0.5) (5.6,0.5)};
\node[anchor=south west, font=\scriptsize, text=black!60] at (axis cs:0.45,0.5) {no signal};
\addplot[only marks, mark=*, mark size=2.3pt, cGrd, error bars/.cd, y dir=minus, y explicit, error mark=none,
  error bar style={black!35, line width=1.2pt}]
  table[x=x, y=fsm, y error expr=\thisrow{fsm}-\thisrow{cstr}] {data/mechanism.dat};
\addlegendentry{path planning (retrieval errors)}
\addplot[only marks, mark=o, mark size=2.3pt, cInt, thick] table[x=x, y=cstr] {data/mechanism.dat};
\addlegendentry{CSTR (reasoning errors)}
\end{axis}
\end{tikzpicture}

\caption{Does low attention on the deciding prompt region predict failure? AUROC of the negative minimum
grounding share $-g_{\min}$ for a failing answer (training and development data; 0.5 = no signal). In path
planning, where errors are retrieval errors, grounding is predictive for every model. In the CSTR, where models
consult the relevant physics and reason incorrectly, it is weak or absent. The CSTR value for Qwen-3B comes from
the exploratory freeze.}
\label{fig:mechanism}
\end{figure}

\paragraph{Answer to RQ3} Attention-based internals detect the first kind of failure and not the second. This
explains the domain dependence of RQ2: internals add most where the agent's errors are failures to consult
information it was given. The finding is consistent with the view that attention is a useful predictive signal
without being a complete explanation of model behaviour \citep{jain2019attention,wiegreffe2019attention}. It
also gives practitioners a test they can run before deployment: compute the grounding AUROC on development
data. If it is near 0.5, internal signals are unlikely to add to a good observable screen.

\subsection{Robustness to plant and fault shift}\label{sec:res-shift}
A screen is only useful if it remains valid when the population of plants or faults changes. In a pre-registered
shift analysis (Table~\ref{tab:shift}) we fitted probes on one part of the plant and fault space and evaluated
them on another, without recalibration. Three shifts were tested: from low to high nominal feed rate (visible in
the prompt), from low to high heat-transfer coefficient (not visible), and each fault family held out in turn.
All signals degraded under the feed and fault-family shifts. The observable screen lost 0.04--0.27 AUROC, whereas
the cooling shift was nearly harmless. Contrary to our pre-registered expectation, internal signals did not
degrade less than observables. For Qwen2.5-7B, the internals-only probe fell from 0.80 to 0.52 under the feed
shift. An earlier and larger shift, from one training plant to a varied population under a different prompt
version, had favoured internal signals. We report both results and conclude that, in our data, internal signals
offer no special robustness. A screen must be recalibrated, and re-certified, when the plant or fault population
changes.

\begin{table}[!htbp]
\centering
\small
\caption{AUROC for failure on the target part of each pre-registered shift (probes fitted on the source part only). Control: random split of the same size.}
\label{tab:shift}
\begin{tabular}{llrrrr}
\toprule
Model & Signal & Control & Feed & Cooling & Fault family \\
\midrule
Qwen-1.5B & Observables & 0.93 & 0.70 & 0.91 & 0.72 \\
 & Obs. + all internals & 0.92 & 0.72 & 0.90 & 0.74 \\
 & All internals & 0.80 & 0.64 & 0.79 & 0.66 \\
 & Token confidence & 0.65 & 0.54 & 0.63 & 0.50 \\
\midrule
Qwen-3B & Observables & 0.87 & 0.73 & 0.86 & 0.80 \\
 & Obs. + all internals & 0.85 & 0.74 & 0.85 & 0.74 \\
 & All internals & 0.71 & 0.60 & 0.72 & 0.62 \\
 & Token confidence & 0.57 & 0.48 & 0.55 & 0.50 \\
\midrule
Qwen-7B & Observables & 0.93 & 0.66 & 0.92 & 0.77 \\
 & Obs. + all internals & 0.92 & 0.65 & 0.91 & 0.71 \\
 & All internals & 0.80 & 0.52 & 0.77 & 0.55 \\
 & Token confidence & 0.68 & 0.47 & 0.66 & 0.54 \\
\midrule
Llama-3B & Observables & 0.90 & 0.86 & 0.91 & 0.80 \\
 & Obs. + all internals & 0.92 & 0.91 & 0.93 & 0.82 \\
 & All internals & 0.84 & 0.83 & 0.86 & 0.71 \\
 & Token confidence & 0.63 & 0.57 & 0.66 & 0.57 \\
\bottomrule
\end{tabular}
\end{table}

